% Praktikumsarbeitsblatt Versuch 2 — Zustandsschätzung (Kalman-Filter) im 2D-Simulator.
% Bauen: make kalman  (latexmk -pdf). Zwei TikZ-Figuren, keine externen Grafiken nötig.
% Stil und Präambel folgen docs/praktikum/anleitung.tex (Versuch 1) — gleiche Pakete.
\documentclass[12pt,a4paper]{scrartcl}
\usepackage[ngerman]{babel}
\usepackage[T1]{fontenc}
\usepackage[utf8]{inputenc}
\usepackage{lmodern}
\usepackage[a4paper,left=2.8cm,right=2.6cm,top=3.0cm,bottom=2.8cm]{geometry}
\usepackage{amsmath,amssymb}
\usepackage{booktabs}
\usepackage{graphicx}
\usepackage[table]{xcolor}
\usepackage{listings}
\usepackage{tikz}
\usetikzlibrary{positioning,calc,arrows.meta}
\usepackage[font=small,skip=4pt]{caption}
\usepackage{enumitem}
\usepackage{microtype}
\usepackage[colorlinks=true,linkcolor=blue!55!black,urlcolor=blue!55!black,
            citecolor=blue!55!black]{hyperref}
\usepackage{cleveref}

\lstset{language=Python, basicstyle=\ttfamily\small, columns=fullflexible,
        keepspaces=true, tabsize=4, showstringspaces=false, breaklines=true,
        keywordstyle=\bfseries, commentstyle=\itshape\color{black!55},
        stringstyle=\color{black!65}, frame=l, framerule=0.3pt,
        rulecolor=\color{black!30}, xleftmargin=6pt, framexleftmargin=6pt,
        aboveskip=8pt, belowskip=6pt, numbers=none, escapechar=|}

\newcommand{\thema}[1]{\texttt{#1}}
\newcommand{\ordner}[1]{\texttt{#1}}
\newcommand{\kommando}[1]{\texttt{#1}}
\newcommand{\xhat}{\hat{x}}
\setlist[enumerate,1]{leftmargin=*,itemsep=2pt}
\setlist[itemize,1]{leftmargin=*,itemsep=1pt,topsep=3pt}
\hyphenation{Messwerte Zustandsschaetzung Covarianz}

\title{\textbf{Zustandsschätzung im Simulator}\\[2pt]
  \large Praktikumsversuch 2 — Kalman-Filter für GPS, Odometrie und IMU}
\author{Betreuung: Robotik-Labor}
\date{Version 0.1 — Wintersemester}

\begin{document}
\maketitle

\begin{abstract}
\noindent
Sie schätzen, Sie fahren nicht. Ein Bewerter setzt Ihren Roboter an den Start, fährt
ihn durch eine vorgegebene Kommandofahrt und zeichnet dabei die exakte Pose
(\thema{truth}) und Ihre Schätzung \thema{kf/pose} auf. Urteilt wird über Genauigkeit:
RMSE gegen die Wahrheit, Verbesserungsverhältnis gegenüber dem Rohsensor, Maximalfehler
während eines GPS-Funklochs und Konsistenz der von Ihnen angegebenen Standardabweichung
(NEES). Dazu bauen Sie einen Kalman-Filter mit Zustandsvektor aus Position und
Geschwindigkeit, füttern ihn mit GPS (langsam, verrauscht, zeitweise weg), Odometrie
(schnell, driftend) und IMU (kurzfristig sauber, langfristig wertlos) und melden
\emph{mit}, wie unsicher Sie sind. Der Simulator ist derselbe aus Versuch 1 — neu sind
ein IMU-Rauschmodell, drei Themen und die Frage, ob Ihre Zahlen stimmen.
\end{abstract}

\tableofcontents

\section{Versuchsziel und Vorfragen}\label{sec:vorfragen}
\subsection{Ziel}
Nach dem Versuch können Sie
\begin{itemize}
  \item einen Zustandsvektor \((x,y,v_x,v_y)\) mit einem CV-Modell prädizieren und jede
        Messung mit ihrer eigenen Unsicherheit einarbeiten — Prädiktion, Update und
        Kalman-Verstärkung aus \(P\), \(Q\) und \(R\), nicht aus dem Bauch,
  \item die eigene Unsicherheit so angeben, dass sie \emph{stimmt}, und sie mit dem
        normierten Fehlerquadratmittel (NEES) prüfen,
  \item Odometrie-\emph{Twist} und Gyro-Drehrate als Bewegungsmodell benutzen, um ein
        GPS-Funkloch von zehn Sekunden unter 1,8\,m Fehler zu überbrücken,
  \item Drift (korreliert, wachsend) von Messrauschen (weiß, mittelwertfrei) trennen und
        begründen, warum eine IMU niemals offen integriert wird,
  \item aus einem Messprotokoll RMSE-Zeitreihen, eine \(Q\)/\(R\)-Tabelle und eine
        NEES-Diskussion für den Bericht erzeugen.
\end{itemize}

\begin{quote}\small
\textbf{Rollenverteilung.}\ \newline
Der Bewerter setzt den Roboter vor \emph{jedem} KF-Auftrag neu an der Spawn-Pose ab und
fährt die Kommandofolge aus \ordner{config/tasks.json} selbst über \thema{cmd\_vel}.
Ihr Knoten sendet \emph{keine} Radgeschwindigkeiten und definiert keine
\thema{inverse\_kinematics} — sonst konkurriert er mit dem Bewerter und der Lauf ist
wertlos. Gearbeitet wird allein pro Person: ein Robotername, ein Filterknoten.
\end{quote}

\subsection{Vorfragen (vor dem Praktikumsstermin rechnen)}
Verwenden Sie die Simulatorwerte aus \cref{tab:parameter}
(\ordner{mecanum\_lab/types.py}: \texttt{DEFAULT\_CONFIG}) und die Prüfprofile aus
\ordner{config/tasks.json}: Rauschendichte \(\sigma_a=0{,}002\,\mathrm{m/s^2/\sqrt{Hz}}\),
Accelerometer-Bias \(b_a=0{,}05\,\mathrm{m/s^2}\), IMU-Rate \(100\,\mathrm{Hz}\),
Radrauschen \(\sigma_\text{wheel}=0{,}04\,\mathrm{rad/s}\), GPS-Streuung
\(0{,}5\,\mathrm{m}\). Lösungen: \cref{sec:loesungen}.

\begin{enumerate}
  \item \textbf{Was dominiert die IMU?} Eine Rauschendichte von
        \(0{,}002\,\mathrm{m/s^2/\sqrt{Hz}}\) bei \(100\,\mathrm{Hz}\): welche
        Standardabweichung hat \emph{eine} Stichprobe, und warum steht dort
        \(\sqrt{\text{rate}/2}\)? Rechnen Sie den Wegfehler aus, den ein fester Bias von
        \(0{,}05\,\mathrm{m/s^2}\) nach 10\,s Doppelintegration liefert, und vergleichen
        Sie ihn mit dem Beitrag des weißen Rauschens nach derselben Zeit. Was folgt
        daraus für die Filterpraxis?
  \item \textbf{Rekursion von Hand.} Eindimensionales CV-Modell mit Zustand
        \((x,v)^{\mathsf T}\), \(F=\left(\begin{smallmatrix}1&dt\\0&1\end{smallmatrix}\right)\),
        \(H=\begin{pmatrix}1&0\end{pmatrix}\). Gegeben: \(\hat x=2{,}0\,\mathrm{m}\),
        \(\hat v=0\), \(P=\operatorname{diag}(0{,}04;\ 1{,}0)\) in \(\mathrm{m^2}\) bzw.
        \(\mathrm{m^2/s^2}\), \(dt=0{,}5\,\mathrm{s}\), \(q=0{,}6\,\mathrm{m^2/s^3}\),
        \(R=0{,}25\,\mathrm{m^2}\), Messung \(z=2{,}40\,\mathrm{m}\). Bestimmen Sie
        \(P^{-}\), die Verstärkung \(K\) (beide Einträge), \(\hat x^{+}\), \(\hat v^{+}\)
        und die Standardabweichung der Position vor und nach dem Update.
  \item \textbf{Drift ist kein weißes Rauschen.} Die Odometrie entsteht aus den
        \emph{gemessenen} Radgeschwindigkeiten, und auf jeder steckt \(\sigma_\text{wheel}\).
        Warum darf der \emph{Pose}-Fehler der Odometrie nicht als weißes Messrauschen mit
        konstantem \(R\) in den Filter? Was machen Kovarianz und NEES, wenn es doch
        getan wird?
  \item \textbf{NEES.} Die Auswertung meldet für eine Fahrt \(\overline{\text{NEES}}=0{,}1\)
        bzw. \(1\) bzw. \(8\). Was heißt das jeweils für die angegebene
        Standardabweichung (Faktor, welche Richtung), und welcher der drei Werte fällt
        durch das Band für K3 (\(0{,}15\ldots3{,}5\))?
\end{enumerate}

\section{Theorie: prädizieren, messen, abwägen}\label{sec:theorie}
\subsection{Zustand, Prädiktion, Update, Kovarianz}
Der Filter kennt einen Zustandsvektor \(\boldsymbol\xi\) und dessen Unsicherheit, die
Kovarianz \(P\). Hier genügt zunächst \(\boldsymbol\xi=(x,y,v_x,v_y)^{\mathsf T}\) in
Weltmetern; K4 verlangt einen weiteren Freiheitsgrad für die Drehrate. Es läuft immer
derselbe Zyklus: \textbf{Prädiktion} über die echte Zeitspanne \(dt\) seit dem letzten
Zustand, \textbf{Update}, sobald eine Messung \(z\) mit ihrer Unsicherheit \(R\) eintrifft.
\begin{align}
  \boldsymbol\xi^{-} &= F\,\boldsymbol\xi + \text{Antriebe},
   && P^{-} = F\,P\,F^{\mathsf T} + Q
   \label{eq:pradiktion}\\
  \mathbf y &= z - H\,\boldsymbol\xi^{-},
   && S = H\,P^{-}H^{\mathsf T} + R,\quad K = P^{-}H^{\mathsf T}S^{-1}
   \label{eq:innovation}\\
  \boldsymbol\xi^{+} &= \boldsymbol\xi^{-} + K\,\mathbf y,
   && P^{+} = (I-KH)P^{-}(I-KH)^{\mathsf T} + KRK^{\mathsf T}
   \label{eq:update}
\end{align}
Bei einer direkt gemessenen Größe ist \(K = P^{-}/(P^{-}+R)\): ist die Prognose
unsicherer als die Messung, zieht der Filter zur Messung; ist \(P\) klein, passiert
wenig. Die Stellschrauben bedeuten Verschiedenes: \(Q\) sagt, was dem \emph{Modell} fehlt
(unerklärte Beschleunigung), \(R\), was dem \emph{Sensor} fehlt. Wer nur \(Q\) erhöht,
macht den Filter unruhig; wer nur \(R\) erhöht, macht ihn träge. Beide bestimmen die
Kovarianz — und damit die Kennzahlen, mit denen bewertet wird:
\begin{equation}
  \mathrm{NEES} = \frac1N\sum_i \frac12\!
     \left(\frac{e_{x,i}^{2}}{s_{x,i}^{2}} + \frac{e_{y,i}^{2}}{s_{y,i}^{2}}\right),\quad
  \mathrm{RMSE} = \sqrt{\langle e_x^2+e_y^2\rangle},\quad
  \text{Verb.} = \frac{\mathrm{RMSE}_{\text{GPS}}}{\mathrm{RMSE}_{\text{KF}}}.
  \label{eq:nees}
\end{equation}
Bei zwei Freiheitsgraden ist der erwartete NEES-Wert 1; ein Einzelwert bis etwa 5,99 ist
bei \(\chi^2\)-Verteilung noch unauffällig. Ohne \cref{eq:update} in der Joseph-Form
bleibt \(P\) nach Rundungsfehlern gern asymmetrisch — bei einem 4×4-Filter von Hand
vermutlich die häufigste stille Quelle für NEES weit über 1.

\subsection{Das CV-Modell in Matrizen}
Konstante Geschwindigkeit heißt das Modell, weil es die Beschleunigung nicht kennt. In
Matrizen (Weltframe, Radreihenfolge irrelevant, Zustand \((x,y,v_x,v_y)\)):
\begin{equation}
  F(dt) = \begin{pmatrix}1&0&dt&0\\0&1&0&dt\\0&0&1&0\\0&0&0&1\end{pmatrix},\quad
  H = \begin{pmatrix}1&0&0&0\\0&1&0&0\end{pmatrix},\quad
  R = \begin{pmatrix}\sigma_{\text{gps}}^2&0\\0&\sigma_{\text{gps}}^2\end{pmatrix},
  \label{eq:cv}
\end{equation}
\begin{equation}
  Q(dt) \;=\; q\,
  \begin{pmatrix}\frac{dt^4}{4}&0&\frac{dt^3}{2}&0\\[2pt]
                0&\frac{dt^4}{4}&0&\frac{dt^3}{2}\\[2pt]
                  \frac{dt^3}{2}&0&dt^2&0\\[2pt]
                0&\frac{dt^3}{2}&0&dt^2\end{pmatrix},
  \qquad q\ \mathrm{in\ m^2/s^3}.
  \label{eq:q}
\end{equation}
\(q\) ist die angenommene Unsicherheit der \emph{Beschleunigung}, also dessen, was das
CV-Modell nicht darstellen kann; brauchbarer Start \(0{,}3\ldots1{,}0\), empfohlen in
\texttt{kf.q\_acc} = 0,6 (\cref{tab:parameter}). \(\sigma_{\text{gps}}\) wird nicht
geraten, sondern aus dem Sensorprofil gelesen
(\kommando{rob.config("gps.sigma\_xy", 0.5)} bzw. \thema{/sim/config}). Ohne numpy reicht
handschriftliche Matrizenmultiplikation: \cref{eq:cv} ist dünn besetzt, \cref{eq:q}
blockweise — vier Zeilen Code je Achse. Zwei Dinge entscheiden über Gut oder Schlecht:
\(dt\) ist die \emph{echte} Differenz der Messstempel (nicht ein gemittelter Fixwert), und
es wird bis zum \emph{Messzeitpunkt} der eintreffenden Messung prädiziert, erst dann
aktualisiert. Bei 5\,Hz GPS und 0,85\,m/s kosten 100\,ms Rückstand 8,5\,cm — das ist
bereits ein Drittel des K4-Ziels.

\subsection{Warum die Gyro kurzfristig und die Odometrie langfristig hilft}
Zwei Fakten bestimmen die Architektur des Filters. \textbf{Erstens:} eine IMU misst
Specific Force, und ihr Fehler wächst mit der Integrationsordnung. Die Werte aus
\cref{tab:parameter} sagen: Gyro-Rauschen \(10^{-4}\,\mathrm{rad/s/\sqrt{Hz}}\) sind pro
Stichprobe \(7{,}1\cdot10^{-4}\,\mathrm{rad/s}\) — vernachlässigbar; der Gyro-Bias
\(5\cdot10^{-3}\,\mathrm{rad/s}\) ergibt in 10\,s immerhin \(0{,}05\,\mathrm{rad}\)
(\(2{,}9^{\circ}\)) Kursfehler, der auf 10\,m Fahrt \(0{,}25\,\mathrm{m}\) Querfehler
kostet. Beim Accelerometer wird aus \(0{,}05\,\mathrm{m/s^2}\) Bias durch
\emph{zwei} Integrationen in 10\,s Wegfehler von \(\tfrac12 b t^2 = 2{,}5\,\mathrm{m}\).
Deshalb: Drehrate ja (einmal integrieren, klein), Beschleunigung nur kurzfristig und nie
offen — die IMU ist ein Bewegungsmodell für Sekundenbruchteile, kein Ortungsgerät.
\textbf{Zweitens:} die Odometrie liefert zwei verschiedene Dinge. Ihre \emph{Pose} ist
Dead Reckoning und driftet (Zufalls-spaziergang aus \(\sigma_\text{wheel}\),
kursiert durch den Drehfehler). Ihr \emph{Twist} \((v_x,v_y,\omega)\) dagegen ist eine
unmittelbare Messung: mit \(r=0{,}05\,\mathrm{m}\) und \(\sigma_\text{wheel}=0{,}04\) bis
\(0{,}05\,\mathrm{rad/s}\) steckt auf \(v_x=\tfrac{r}{4}\sum_i w_i\) nur
\(\approx 1\,\mathrm{mm/s}\) Rauschen \emph{ohne} Drift. Der gangbarste Filter tut daher:
Prädiktion mit der \emph{gemessenen} Körpergeschwindigkeit (Odometrie-Twist, in die Welt
gedreht) und der Gyro-Drehrate, Update der Position mit GPS. Wer stattdessen die
Odometrie-\emph{Pose} als Messung mit festem \(R\) einwirft, doppelt information, die schon
in der Prädiktion steckt — der dritte Weg ist, die Drift als eigenen Zustandswert
(Bias) mitzuführen und \(R\) mit der Zeit wachsen zu lassen.

\subsection{GPS: Rauschen, Bias, Funkloch, Modellfehler}
Das GPS ist der einzige absolute Sensor. Sein Rauschen ist unbekannt korreliert, also
geradezu gemacht für \(R\): \(\sigma_{xy}=0{,}5\,\mathrm{m}\) in K1/K3,
\(0{,}8\,\mathrm{m}\) in K2, \(0{,}6\,\mathrm{m}\) in K4 — die Roh-RMSE liegt bei
\(\sigma\sqrt2\approx0{,}7\,\mathrm{m}\) bzw. \(1{,}1\,\mathrm{m}\), und das
Verbesserungsverhältnis in \cref{eq:nees} wird gegen genau diese Zahl gerechnet. Ein
\emph{fester} Bias (\texttt{gps.bias\_xy}, Multi-Path) ist aus GPS allein
unbeobachtbar: kein Filter kann ihn entfernen, wohl aber als konstante Unsicherheit in
\(R\) tragen; \texttt{gps.bias\_step} legt ihn für ein Zeitfenster obendrauf und ist die
Ausreißer-Übung. Das \textbf{Funkloch} schließlich (\texttt{gps.gap = [17, 10]} in K2)
entwässert die Messseite vollständig: Die Prädiktion läuft weiter, die Kovarianz wächst
nur noch über \(Q\), und der Fehler wächst mit der verstrichenen Zeit — begrenzt durch den
Modellfehler, denn ein CV-Modell kann eine Kurve mit \(0{,}26\,\mathrm{rad/s}\) nicht
vorhersagen. Deshalb steht in K2 nicht die Mittel-RMSE im Vordergrund, sondern der
Maximalfehler in der Lücke (\(\le 1{,}8\,\mathrm{m}\)): Wer die Drehrate aus dem Gyro in
den Zustandsvektor nimmt, übersteht die zehn Sekunden; wer rein CV betreibt, läuft
bogenau der Hallenecke hinterher.

\begin{figure}[ht]
  \centering
  \begin{tikzpicture}[
    every node/.style={font=\small},
    blk/.style={draw=black!70,line width=0.6pt,rounded corners=2pt,fill=black!5,
                align=center,inner sep=4pt,minimum height=9mm},
    sim/.style={blk,fill=blue!7,draw=blue!50!black},
    stu/.style={blk,fill=green!9,draw=green!45!black,very thick},
    bw/.style={blk,fill=orange!16,draw=black!60},
    fl/.style={-{Stealth[length=2mm]},line width=0.7pt,draw=black!70},
    fld/.style={-{Stealth[length=2mm]},line width=0.6pt,draw=black!45,dashed},
    lbl/.style={font=\scriptsize\ttfamily,inner sep=1.5pt,fill=white,fill opacity=0.85},
  ]
  \node[bw] (gr) at (0,0) {Bewerter \\ \scriptsize grade.Grader};
  \node[sim,right=20mm of gr] (eng) {Simulator \\ \scriptsize engine + sensors,
    50\,Hz Physik};
  \node[stu,below=22mm of eng] (kf) {Ihr Filter \\ \scriptsize student/kf\_template.py,
    mission(rob, task)};
  \node[blk,below=22mm of gr] (gui) {GUI \scriptsize\ render.py};
  \draw[fl] (gr) -- node[lbl,above]{$\langle$robot$\rangle$/cmd\_vel} (eng);
  \draw[fld] (eng) -- node[lbl,above,yshift=1pt]{$\langle$robot$\rangle$/truth
    (nur debug\_truth)} (gr);
  \draw[fl] ([xshift=9mm]eng.south) -- node[lbl,right,align=left] {$\langle$robot$\rangle$/imu
    100\,Hz \\ $\langle$robot$\rangle$/gps 1\,bis 5\,Hz \\ $\langle$robot$\rangle$/odom 50\,Hz
    \\ /sim/config} ([xshift=9mm]kf.north);
  \draw[fl] ([xshift=-9mm]kf.north) -- node[lbl,left,align=right] {$\langle$robot$\rangle$/kf/pose
    \\ \itshape kf/info} ([xshift=-9mm]eng.south);
  \draw[fld] (gui) -- node[lbl,below]{$k$: Spur, GPS, Sch\"atzung, $2\sigma$} (kf);
  \draw[fld] (gui.east) -- node[lbl,above]{$\langle$robot$\rangle$/odom, kf/pose} (eng.south west);
  \node[font=\scriptsize,align=left,anchor=north west] at ($(gui.south west)+(0,-9mm)$)
    {\emph{nie} von Ihnen gesendet: \ttfamily wheel\_speeds, cmd\_vel};
  \end{tikzpicture}
  \caption{Signalfuss von Versuch 2. Der Bewerter fährt, die Simulation liefert die
    Sensoren, Sie liefern \thema{kf/pose} \emph{inklusive} Standardabweichung —
    an der gemessenen Stelle, nicht an der Wahrheit.}
  \label{fig:signalfuss}
\end{figure}

\begin{figure}[ht]
  \centering
  \begin{tikzpicture}[x=0.40cm, every node/.style={font=\scriptsize\sffamily}]
  \draw[-{Stealth[length=1.8mm]}] (-0.4,0) -- (39.2,0);
  \foreach \t in {0,5,...,35} {\draw[black!60] (\t,1.4mm) -- (\t,-1.4mm)
    node[below=0.5pt,black!70] {\t};}
  \node[anchor=west] at (37.6,-0.6) {$t$ in s};
  \draw[fill=blue!14,draw=blue!55!black] (0,0.45) rectangle (6,1.15)
    node[midway,font=\scriptsize\sffamily] {$v_x=0{,}6$};
  \draw[fill=blue!26,draw=blue!55!black] (6,0.45) rectangle (31,1.15)
    node[midway] {$v_x=0{,}6$, $\omega=0{,}26$ + überlagerte Sinusschwingung};
  \draw[fill=blue!14,draw=blue!55!black] (31,0.45) rectangle (37,1.15)
    node[midway] {$v_x=-0{,}6$};
  \node[align=right,anchor=south east] at (0,1.25) {Kommandofahrt};
  \foreach \t in {0.5,1.5,...,16.5,27.5,28.5,...,37.5}
    \draw[black!75] (\t,1.75) -- (\t,2.25);
  \node[anchor=south east] at (0,2.3) {GPS-Fixe 1\,Hz};
  \draw[fill=red!22,draw=red!60!black] (17,1.65) rectangle (27,2.35)
    node[midway,text=black] {Funkloch \ttfamily gps.gap=[17,10]};
  \draw[dashed,black!70] (4,-0.9) -- (4,3.1)
    node[above,align=center] {\ttfamily einlauf: 4\\ \scriptsize hier beginnt die Messung};
  \draw[fill=green!20,draw=green!45!black] (0,-0.75) rectangle (37,-0.3)
    node[midway] {Ihre Meldungen auf $\langle$robot$\rangle$/kf/pose, mindestens 5\,Hz};
  \node[anchor=north west,text=red!60!black] at (17,-0.95)
    {hier zählt der Maximalfehler: $\le 1{,}8$\,m};
  \end{tikzpicture}
  \caption{Zeitstrahl des Auftrags \thema{kf\_fusion}: Kommandosegmente aus
    \ordner{config/tasks.json}, GPS-Fixe mit Lücke von 17\,s bis 27\,s, und der
    Messbereich des Bewerters ab \thema{einlauf}.}
  \label{fig:zeitstrahl}
\end{figure}

\section{Der Simulator}\label{sec:simulator}
\subsection{Neue Themen und was sie enthalten}
Versuch 1 bleibt, wie er ist; der Filter setzt darauf. \cref{tab:themen} zeigt den Zugang
— die Pfade hängen am Roboternamen, \kommando{topic("kf","alice")} ergibt
\thema{/alice/kf/pose}.

\begin{table}[ht]\centering\small
\caption{Neue und geänderte Themen in Versuch 2 (\ordner{types.MSG\_SPECS}).}
\label{tab:themen}
\begin{tabular}{@{}p{4.1cm}p{5.4cm}p{5.5cm}@{}}
\toprule
Thema & Typ (ROS 2) & Inhalt und Richtung \\
\midrule
\thema{<robot>/imu} & \thema{sensor\_msgs/msg/Imu}
  & 100 Hz (einstellbar), Achsen im Körperframe, \thema{linear\_acceleration.z}
  \(\approx +g\); \thema{roll}/\thema{pitch} mitgeliefert \\
\thema{<robot>/kf/pose} & \thema{geometry\_msgs/msg/PoseWith\_Covariance\_Stamped}
  & Ihre Schätzung inklusive 1\(\sigma\) (Diagonale) — die Bewertungsgröße \\
\thema{<robot>/kf/info} & \thema{std\_msgs/msg/String} (JSON)
  & Diagnose (Laufzeit, Raten, \(Q\)/\(R\)); der Auswerter ignoriert sie \\
\thema{<robot>/truth} & \thema{geometry\_msgs/msg/PoseStamped}
  & exakte Pose mit einstellbarer \thema{truth.rate}, nur mit \thema{debug\_truth}
  (= \kommando{--truth}) \\
\thema{/sim/config} & \thema{std\_msgs/msg/String} (JSON)
  & aktives Sensorprofil (\thema{gps}, \thema{odom}, \thema{imu}, \thema{truth},
  \thema{rate}, \thema{seed}) — hier steht \(\sigma\), das Sie in \(R\) einsetzen \\
\thema{<robot>/odom}, \thema{<robot>/gps} & wie Versuch 1
  & Odometrie 50 Hz mit \thema{twist}; GPS mit Rate und Streuung des Prüfprofils \\
\bottomrule
\end{tabular}
\end{table}

\subsection{Das IMU-Modell und seine Zahlen}
\ordner{mecanum\_lab/types.py} (\texttt{DEFAULT\_CONFIG}) ist die Quelle aller Zahlen;
\ordner{config/default.json} und \kommando{--set} darüber schreiben, das Prüfprofil des
Auftrags sowieso (\cref{sec:inbetriebnahme}). Die Spalte „pro Stichprobe'' ist die für
den Filter relevante Größe: Dichten werden mit \(\sqrt{\text{rate}/2}\) auf die Samplingrate
umgerechnet, Bias-Random-Walk mit \(\sqrt{1/\text{rate}}\) pro Schritt.

\begin{table}[ht]\centering\small
\caption{Sensorik-Einstellungen aus \texttt{DEFAULT\_CONFIG} (Versuch-2-Auszug).}
\label{tab:parameter}
\begin{tabular}{@{}llp{7.6cm}@{}}
\toprule
Größe & Wert & Bedeutung für Ihren Filter \\
\midrule
\multicolumn{3}{@{}l}{\textit{IMU (\texttt{imu}) — MEMS-Klasse, absichtlich so kaputt}}\\
\thema{rate} & 100 Hz & K4 fährt 200 Hz; \kommando{--set imu.rate=400} ist erlaubt \\
\thema{accel\_noise} & 2\,e-3 m/s²/\(\sqrt{\mathrm{Hz}}\) & \(\Rightarrow\) 14 mm/s² pro Stichprobe \\
\thema{accel\_bias} & 0,05 m/s² & je Roboter und Achse neu gezogen (Seed); doppelte
  Integration: 2,5 m in 10 s \\
\thema{accel\_bias\_walk} & 5\,e-4 m/s²/\(\sqrt{\mathrm{s}}\) & der Bias bleibt nicht, wo er war: nach 10 s
  streut er mit 1,6\,mm/s² \\
\thema{accel\_scale} & 1\,e-3 (0,1\,\%) & verschwindet nicht durch Mitteln \\
\thema{gyro\_noise} & 1\,e-4 rad/s/\(\sqrt{\mathrm{Hz}}\) & \(\Rightarrow\) 0,7\,mrad/s pro Stichprobe \\
\thema{gyro\_bias} & 5\,e-3 rad/s & 0,05 rad Kursfehler in 10 s — mitteln hilft \\
\thema{gyro\_bias\_walk} & 2\,e-5 rad/s/\(\sqrt{\mathrm{s}}\) & nach 10 s streut der Bias selbst
  mit 0,06\,mrad/s — klein, aber er verschwindet nicht von allein \\
\thema{tilt\_sigma}, \thema{tilt\_tau} & 6\,mrad, 0,4 s & Federung wackelt, Schwerkraft
  sickert in \thema{ax}/\thema{ay}: \(g\varphi \approx 0{,}06\) m/s² \\
\thema{vibration}, \thema{vibration\_hz} & 0,08 m/s², 16 Hz & Fahrwerksschwingung,
  mittelbar (Tiefpass oder Mittelung) \\
\thema{startup}, \thema{startup\_bias} & 0,4 s, 0,5 m/s² & Einschwingen des Treibers —
  in dieser Zeit nicht kalibrieren \\
\thema{gravity} & 9,81 m/s² & steckt voll in \thema{az} \\
\addlinespace
\multicolumn{3}{@{}l}{\textit{Übrige Sensorik}}\\
\thema{gps.rate}, \thema{gps.sigma\_xy} & 5 Hz, 0,06 m & Versuch 1; die KF-Aufträge
  setzen 1\,bis 5 Hz und 0,5\,bis 0,8 m \\
\thema{gps.bias\_xy}, \thema{gap}, \thema{bias\_step} & \([0,0]\), null, null
  & Versatz, Funkloch \([t_0,\text{dauer}]\), springender Bias \\
\thema{odom.rate}, \thema{sigma\_wheel} & 50 Hz, 0,04 rad/s & auf jeder Radmessung —
  daraus wird die Drift \\
\thema{truth.rate} & 20 Hz & Rate der exakten Pose (Bewertung) \\
\thema{kf.rate}, \thema{kf.q\_acc}, \thema{kf.q\_turn} & 20 Hz, 0,6 m²/s³, 0,02
  & \emph{Empfehlung an Sie} — die Simulation benutzt diesen Block nicht \\
\bottomrule
\end{tabular}
\end{table}

\textbf{Warum \thema{az} im Stand \(\approx +9{,}81\) anzeigt.} Ein
MEMS-Beschleunigungsmesser misst keine Bewegung, sondern die Specific Force — die Kraft
pro Masse, die ihn am freien Fall hindert. Auf ebener Fläche ist das die Normalkraft nach
oben, also \(+g\). Das ist kein Messfehler und kein Vorzeichenbug, sondern die Physik des
Chips: wer \thema{az} ohne Weiteres integriert, bekommt in 10\,s 490\,m Höhe. Für die
Ebene zählen nur \thema{ax} und \thema{ay} — und die sind verbiasst, geneigt und
überlagert, eben deshalb nur kurzfristig brauchbar. \thema{gx}/\thema{gy} sind durch
dieselbe Neigung mitbedingt, \thema{gz} ist das Hauptsignal in 2D.

\textbf{Nachmessen statt glauben.} Ein Standlauf über 5\,s (Seed 1, Roboter \emph{alice},
502 Stichproben) liefert \thema{az} \(=9{,}74\pm0{,}014\,\mathrm{m/s^2}\) — die Streuung
ist exakt die gerechnete Dichte, der Unterschied zu 9,81 ist der \thema{accel\_bias} der
\(z\)-Achse — und ein Gyro-Mittel von \(-5{,}4\cdot10^{-3}\,\mathrm{rad/s}\). Genau
dieses Mittel ist Ihre Kalibrierung: IMU im Stillstand mitteln (vor der Fahrt oder in
jederr Pause), Mittel abziehen — so steht die Kalibrierung später auf dem echten Roboter.

\subsection{Auswertung in der GUI}
Die Anzeige ist dieselbe wie in Versuch 1, neu ist die Überlagerung der Schätzung.
Die Hilfeliste im HUD: \kommando{q Ende | SPACE Pause | l Scan | t Spur | k Schätzung |
0 alles | 1..9 Roboter | +/- Zoom}. Die Taste \emph{k} schaltet GPS-Kreuze, die
Schätzung und die \(2\sigma\)-Ellipse drumherum — die Ellipse ist Ihr NEES mit Augen:
sie muss die Wahrheit \emph{umfahren}, nicht umschließen-wollen. Im HUD steht
darüber der aktuelle Abstand Ihrer Schätzung zur Wahrheit (\texttt{kf\_err}) — nur
sichtbar, nicht im Filter benutzt. Für den Bericht gehört die Messreihe in eine Datei:
\kommando{--log messung.csv} schreibt alle 0,05\,s Simulationszeit eine
Semikolon-Zeile mit \thema{t}, \thema{x\_wahr}, \thema{x\_gps}, \thema{t\_alt\_gps},
\thema{x\_kf}, \thema{sx\_kf}, \thema{ax\_imu} und mehr (Spaltenkopf aus
\ordner{mecanum\_lab/logbook.py}); \kommando{python3 tools/kfplot.py messung.csv} malt
Zeitreihen, RMSE und — ohne matplotlib — ein ASCII-Diagramm.

\section{Inbetriebnahme}\label{sec:inbetriebnahme}
\subsection{Installation}
\begin{lstlisting}[language=bash]
./install.sh            # einmalig: prueft pygame, ROS 2 und die Tests
# alternativ:  sudo apt install python3-pygame     (reicht fuer den Stub-Modus)
./lab docs              # Themen, Auftraege und Beispiele auf einen Blick
\end{lstlisting}
Getestet ist Ubuntu 24.04 mit ROS 2 Kilted/Jazzy; ohne ROS läuft alles über den
In-Prozess-Bus (\kommando{MECANUM\_ROS=stub}) — dann zeigt \kommando{ros2 topic list}
natürlich nichts.

\subsection{Werkstatt: den eigenen Filter fahren}
\begin{lstlisting}[language=bash]
# 1) Werkstatt mit Fenster: die Sim faehrt, du schaetzt, truth ist an
./lab run --world arena --task kf_gps --robot alice --robots alice \\
    --controller student/kf_template.py --truth --log messung.csv

# 2) Bewerten wie der Abgabepuefer (ein Auftrag, oder alle vier)
./lab grade --task kf_gps --controller student/kf_template.py --log messung.csv
./lab grade --task kf_alle --controller student/kf_template.py \\
    --log messung.csv --json bericht.json
python3 tools/kfplot.py messung.csv

# 3) Sensorik frei verstellen, keine JSON-Datei noetig
./lab grade --task kf_fusion --controller student/kf_template.py \\
    --set gps.sigma_xy=1.2 --set imu.rate=400 --set gps.gap='[14,8]'
\end{lstlisting}
Drei Fallen beim ersten Start. \emph{Erstens:} \kommando{./lab run} spawn only Roboter aus
\kommando{--robots} — ohne \kommando{--robots alice} läuft eine leere Halle, und das
Protokoll bleibt leer. \emph{Zweitens:} ohne \kommando{--truth} (bzw.
\thema{debug\_truth}) gibt es kein \thema{truth}-Thema, der Bewerter kann dann nicht
messen. \emph{Drittens:} der eigene Knoten darf \thema{wheel\_speeds} und
\thema{cmd\_vel} nicht anfassen (\cref{sec:vorfragen}). Unter ROS gilt außerdem:
\begin{lstlisting}[language=bash]
ros2 topic hz /alice/imu                 # ~100 Hz
ros2 topic echo /alice/imu --once        # az im Stand: rund +9,81
ros2 topic echo /sim/config --once       # sigma_xy, gap, imu-rate des laufenden Profils
ros2 topic hz /alice/kf/pose             # Ihre Rate -- K4 verlangt mindestens 10 Hz
\end{lstlisting}

\subsection{ROS-2-Betrieb mit dem Launch-File}
Ohne Launch-File starten Sie die Simulation mit \kommando{./lab sim}; mit ihm kommt alles
aus einer Zeile — und \emph{jedes} Argument ist eine \kommando{--set}-Angabe an die
Simulator-Konfiguration (\cref{tab:launch}). So verändern Sie Sensorwerte für den Bericht
systematisch, ohne JSON-Dateien anzufassen.

\begin{lstlisting}[language=bash]
source /opt/ros/kilted/setup.bash          # falls ROS nicht automatisch gefunden wird
ros2 launch launch/kf.launch.py --show-args               # alle Argumente erklaeren
ros2 launch launch/kf.launch.py headless:=true aufgabe:=kf_fusion \\
    controller:=student/kf_template.py protokoll:=messung.csv
ros2 launch launch/kf.launch.py bewerten:=kf_alle \\
    controller:=student/kf_solution.py "gps_luecke:=[10,14]" imu_rate:=200
ros2 launch launch/kf.launch.py rviz:=true                # Ansicht mit Kovarianzellipse
\end{lstlisting}

\begin{table}[ht]\centering\small
\caption{Die Argumente, die im Versuch gestellt werden (Auswahl; vollständige Liste:
\kommando{ros2 launch launch/kf.launch.py --show-args}).}
\label{tab:launch}
\begin{tabular}{@{}lp{6.2cm}l@{}}
\toprule
Argument & Wirkung & Standard \\
\midrule
\kommando{aufgabe} & \kommando{kf\_gps}$\vert$\kommando{kf\_fusion}$\vert$\kommando{kf\_kovarianz}$\vert$\kommando{kf\_dynamik} & kf\_gps \\
\kommando{controller} & Ihr Filterknoten, relativ zum Quellbaum & student/kf\_template.py \\
\kommando{bewerten} & Auftrag oder Gruppe (z. B. kf\_alle); leer = nicht bewerten & — \\
\kommando{protokoll} & CSV-Messprotokoll für \ordner{tools/kfplot.py} & — \\
\kommando{gps\_rate}, \kommando{gps\_sigma} & Messungen pro Sekunde, Streuung je Achse in m & Profil \\
\kommando{gps\_luecke} & Funkloch als \kommando{[start, dauer]} in s nach Auftragsbeginn & keins \\
\kommando{gps\_bias\_step} & springender Bias \kommando{[start, dauer, dx, dy]} in m & keiner \\
\kommando{imu\_rate} & IMU-Stichproben pro Sekunde & 100 (K4: 200) \\
\kommando{imu\_accel\_bias}, \kommando{imu\_accel\_walk} & Bias \(1\sigma\) in m/s², Bias-Random-Walk in m/s²/\(\sqrt{s}\) & 0,05 · 5e-4 \\
\kommando{imu\_neigung}, \kommando{imu\_vibration} & Neigung \(1\sigma\) in rad (Schwerkraft sickert in \(a_x\)), Schwingung in m/s² & 0,006 · 0,08 \\
\kommando{wahrheit}, \kommando{seed}, \kommando{sekunden} & truth-Thema, Rausch-Seed, Abbruch nach s & true · 1 · 0 \\
\bottomrule
\end{tabular}
\end{table}

Ein Detail mit Folgen: \kommando{gps\_luecke} und \kommando{gps\_bias\_step} stehen in
Sekunden \emph{nach Auftragsbeginn}, nicht nach Programmstart. Der Bewerter kann das
Funkloch also jederzeit auslösen, und Ihre Messung bleibt dieselbe.

\subsection{Drei Regeln, die den Lauf retten}
\textbf{Erstens: Messstempel statt Wanduhr.} Jede Messung trägt \thema{t} — die
Simulationszeit, in der sie entstanden ist. Benutzen Sie diese Stempel für \(dt\), für
\(Q\) und dafür, ob eine Messung schon verarbeitet ist. Wer stattdessen
\kommando{time.time()} aufruft, bekommt in der Bewertung einen falschen \(dt\): das
Beschleunigungsskript \ordner{tools/fastgrade.py} rechnet 25 Simulationssekunden in einer
Sekunde Wanduhr, und Ihr Filter hält dann eine 5-Sekunden-Lücke für 200 Millisekunden.
\emph{Zweitens: Sie schätzen, Sie fahren nicht.} In den KF-Aufträgen setzt der Bewerter
\thema{cmd\_vel}, die Simulation fährt die Kommandofolge aus \ordner{config/tasks.json}
blind. Ihr Knoten liest Sensoren und antwortet auf \thema{kf/pose} —
\kommando{send\_wheels()}, \kommando{publish\_cmd\_vel()} und \kommando{rob.truth()} sind
in diesem Versuch tabu (der Bewerter merkt das an der fehlenden Fahrt und an der
Kontaktzahl), sonst urteilt der Bewerter über eine Fahrt, die Sie nie gefahren sind.

\emph{Drittens: ein neuer Auftrag beginnt wirklich von vorn.} Wenn sich \kommando{rob.task()}
ändert, wird der Roboter an der Spawn-Pose abgesetzt und das Sensorprofil gewechselt. Beenden
Sie dann Ihren Mess-Loop, lesen Sie \kommando{rob.sensor\_profil()} erneut vor und — das ist
der Fall, der am meisten Punkte kostet — verwerfen Sie die \emph{letzte} Messung des
vorigen Auftrags: Der Bus merkt sich die letzte Meldung, und ein GPS-Fix von vor dem Respawn
ist für den neuen Auftrag eine Lüge. Prüfen Sie also \thema{fix.t} gegen den Beginn Ihrer
Mission (die Musterlösung führt dafür \thema{KF.start}). Die Gyro-Bias-Kalibrierung dürfen Sie
ruhig wiederholen — der Zustand und \(P\) gehören dagegen nicht auf null zurückgesetzt, sonst
werfen Sie Information weg, die Sie gerade erst gesammelt haben.

\section{Aufgaben K1 bis K4}\label{sec:aufgaben}
Alle vier Aufträge fahren dieselbe Logik: Der Bewerter setzt Ihren Roboter an die
Spawn-Pose, fährt die in \ordner{config/tasks.json} hinterlegte Kommandofolge, zeichnet
\thema{truth} und \thema{kf/pose} auf und vergleicht. Gemessen wird erst nach der
Einlaufzeit (4\,s), damit Ihre Bias-Kalibrierung im Stand zählt. \cref{tab:aufgaben} listet
die Prüfgrenzen; \kommando{./lab grade --task kf\_alle --controller student/kf\_template.py}
zeigt Ihnen vorab dieselben Zahlen, die in der Abgabe zählen.

\begin{table}[ht]\centering\small
\caption{Prüfgrenzen der vier Teilaufträge (Quelle: \ordner{config/tasks.json}).}
\label{tab:aufgaben}
\begin{tabular}{@{}lp{3.3cm}p{6.6cm}@{}}
\toprule
Auftrag (Punkte) & Sensorprofil & Prüfgrenzen \\
\midrule
K1 \textit{kf\_gps} (30) & GPS 5 Hz, \(\sigma=0{,}5\) m & RMSE \(\le 0{,}42\) m,
  Verbesserung \(\ge 1{,}6\), Maxfehler \(\le 1{,}25\) m, Rate \(\ge 5\) Hz \\
K2 \textit{kf\_fusion} (30) & GPS 1 Hz, \(\sigma=0{,}8\) m, Funkloch 15…23\,s; Odometrie, IMU
  & RMSE \(\le 0{,}80\) m, Verbesserung \(\ge 2{,}5\), Fehler im Funkloch \(\le 1{,}8\) m \\
K3 \textit{kf\_kovarianz} (20) & wie K1 & mittleres NEES in \([0{,}15\ldots3{,}5]\),
  RMSE \(\le 0{,}42\) m \\
K4 \textit{kf\_dynamik} (10, Bonus) & GPS 5 Hz, \(\sigma=0{,}6\) m, IMU 200 Hz,
  0,85\,m/s mit Schwingung & RMSE \(\le 0{,}35\) m, Verbesserung \(\ge 1{,}8\),
  Maxfehler \(\le 0{,}9\) m, Rate \(\ge 10\) Hz \\
\bottomrule
\end{tabular}
\end{table}

\begin{description}
  \item[K1 — Positionsschätzung aus GPS allein.] Bauen Sie das CV-Modell
    \(\boldsymbol\xi=(x,y,v_x,v_y)^{\mathsf T}\) mit \(F(dt)\) und dem über \(dt\)
    integrierten Beschleunigungsrauschen \(Q\) (\cref{eq:cv,eq:q}) und aktualisieren Sie mit
    der GPS-Position \emph{zum Stempel} \thema{fix.t}. Ziel ist das Verständnis der
    Rekursion, nicht der Trickkiste: ein einzelner, ehrlicher KF reicht.
  \item[K2 — GPS, Odometrie und IMU verschmelzen.] Bei 1 Hz GPS mit \(0{,}8\) m Streuung ist
    der \emph{absolute} Fehler rauschbegrenzt — bewertet wird deshalb die Verbesserung gegen
    das Rohsensorignal. Jetzt hilft der Odometrie-\thema{twist}
    als Eingangsgröße der Prädiktion (in die Welt gedreht) und der Gyro für die Drehrate —
    beides mit Bias, der erst im Stand gemessen werden will. Im Funkloch bleibt nur das
    Bewegungsmodell: Hier entscheidet sich, ob Sie \(Q\) physikalisch begründen oder raten.
  \item[K3 — Wissen, wie gut man schätzt.] Dieselbe Fahrt wie K1, aber bewertet wird die
    \emph{angegebene} Unsicherheit: \thema{kf/pose} trägt \(\sigma_x,\sigma_y\) aus \(P\).
    Vergrößern Sie \(Q\), wenn das NEES über 3 läuft; verkleinern Sie es, wenn es unter 0,3
    fällt — und schreiben Sie den Zielkonflikt in den Bericht.
  \item[K4 — Schnelle Fahrt, treue Schätzung.] 0,85\,m/s mit einer Sinusüberlagerung in der
    Drehrate, IMU mit 200\,Hz. Hier zahlt sich aus, dass Ihr Filter mit der
    \emph{Messrate} läuft statt mit einer Rateschleife.
\end{description}

Typische Fehlerbilder, an denen Sie Ihre Implementierung erkennen: Die Schätzung läuft der
Messung konstant hinterher (\(F\) ohne \(dt\), \(Q=0\)); die Kovarianz schrumpft und
steigt nie wieder (Joseph-Form oder \(P\)-Aktualisierung vergessen); der Roboter wird über
den Gyro gedreht, obwohl der Kurs in der Welt gebraucht wird (Vorzeichen/Drehung der
Körpergeschwindigkeit); das Funkloch wird mit einem \(R\) von \(10^{-6}\) überbrückt (dann
folgt die Schätzung der Odometrie blind und \cref{tab:aufgaben} K2 fällt durch).

\section{Auswertung und Bericht}\label{sec:bericht}
Der Bericht gehört Ihnen — aber die Auswertung liefert Zahlen, die Sie erklären müssen. Er
erwartet:
\begin{itemize}
  \item \textbf{Bahn und Fehlerbild.} Ein Plot aus \ordner{tools/kfplot.py}: Wahrheit, rohes
    GPS, Odometrie, Ihre Schätzung — darunter der Fehler gegen die Wahrheit über der Zeit.
    Beschreiben Sie, \emph{wo} der Fehler wächst (Kurven? Funkloch? nach dem Aufgabenwechsel?).
  \item \textbf{Eine \(Q\)/\(R\)-Tabelle.} Welche Werte haben Sie benutzt, wie sind Sie
    darauf gekommen (Dichte \(\to\) \(Q\), \thema{gps.sigma\_xy} \(\to\) \(R\)), und was
    passiert bei Faktor 3 in beide Richtungen? Vier Messpunkte genügen — aber es müssen
    Messpunkte sein, keine Vermutungen.
  \item \textbf{NEES-Diskussion zu K3.} Ihr Mittelwert, die Streuung innerhalb der Fahrt,
    und was Sie geändert haben, um ins Band zu kommen.
  \item \textbf{Gemessene IMU-Bias-Drift.} Mitteln Sie \thema{imu} im Stand über 10\,s vor
    und nach der Fahrt, geben Sie die Differenz pro Achse an und vergleichen Sie mit
    \cref{tab:parameter}. Daraus folgt Ihre Begründung, warum die IMU nicht offen
    integriert wird.
  \item \textbf{Nutzen der Odometrie im Funkloch.} RMSE im Loch mit und ohne
    Odometrie-Prädiktion (zwei Protokolle, ein Satz Fazit). Wenn Sie behaupten, das Modell
    reiche auch ohne — messen Sie es.
  \item \textbf{Ein Satz zur Modellwahl.} K4 fährt schneller und schwingt: Welche
    Zustandsgröße hätten Sie noch gebraucht (Drehrate im Zustand? Bias als Zustand?), und
    was hätte es gekostet?
\end{itemize}
Rechnen Sie mit dem Protokoll, nicht mit dem Fenster: \kommando{--log messung.csv}
schreibt truth, gps, odom, kf, imu und die Geschwindigkeiten in eine CSV-Datei, die auch
ohne matplotlib auswertbar ist (\kommando{-\-list} zeigt die Spalten).

\appendix
\section{Formelsammlung}\label{sec:formeln}
\textbf{CV-Modell ( geradlinig, Weltframe):}
\[
  F(dt)=\begin{pmatrix}1&0&dt&0\\0&1&0&dt\\0&0&1&0\\0&0&0&1\end{pmatrix}\qquad
  Q(dt)=q\!\begin{pmatrix}\tfrac{dt^3}{3}&0&\tfrac{dt^2}{2}&0\\[2pt]
        0&\tfrac{dt^3}{3}&0&\tfrac{dt^2}{2}\\[2pt]
        \tfrac{dt^2}{2}&0&dt&0\\[2pt]0&\tfrac{dt^2}{2}&0&dt\end{pmatrix}
\]
mit der Beschleunigungs-Dichte \(q\) in m²/s³ — der Simulator startet mit
\(q=0{,}6\) (\thema{kf.q\_acc}), die Drehrate mit \(q_{\text{gier}}=0{,}02\).

\textbf{Rekursion:}
\(\boldsymbol\xi^{-}=F\boldsymbol\xi\), \(P^{-}=FPF^{\mathsf T}+Q\),
\(y=z-H\boldsymbol\xi^{-}\), \(S=HP^{-}H^{\mathsf T}+R\),
\(K=P^{-}H^{\mathsf T}S^{-1}\), \(\boldsymbol\xi=\boldsymbol\xi^{-}+Ky\),
\(P=(I-KH)P^{-}\). Für die Positions-Messung ist \(H\) die Auswahlmatrix mit den beiden
Einsen in der \(x\)- und \(y\)-Spalte, \(R=\operatorname{diag}(\sigma_{\text{gps}}^2,
\sigma_{\text{gps}}^2)\) mit \(\sigma_{\text{gps}}\) aus \thema{gps.sigma\_xy}.

\textbf{Umrechnungen, die man nachschlagen darf:}
Weiße Dichte \(d\) [Einheit/\(\sqrt{\text{Hz}}\)] \(\to\) Streuung pro Stichprobe
\(d\sqrt{\text{rate}/2}\). Bias-Random-Walk \(w\) \(\to\) Streuung pro Schritt
\(w/\sqrt{\text{rate}}\), nach \(T\) Sekunden \(w\sqrt{T}\). Fester Bias \(b\) über \(T\)
Sekunden doppelt integriert: Wegfehler \(\tfrac{1}{2}bT^2\).

\textbf{Konsistenz:}
\(\overline{\text{NEES}}=\frac{1}{N}\sum_i\big[(e_{x,i}^2/\sigma_{x,i}^2
  + e_{y,i}^2/\sigma_{y,i}^2)\big/2\big]\);
erwarteter Wert 1, ein Einzelwert bis 5,99 ist bei zwei Freiheitsgraden noch unauffällig
(\(\chi^2\), 95\,\%). Angegeben werden \(\sigma_x=\sqrt{P_{11}}\),
\(\sigma_y=\sqrt{P_{22}}\) — nicht die Wurzel der Spur.

\section{Lösungen der Vorfragen}\label{sec:loesungen}
\begin{enumerate}
  \item \textbf{IMU.} Eine Stichprobe: \(\sigma_a=0{,}002\cdot\sqrt{100/2}
        = 0{,}014\,\mathrm{m/s^2}\) (14\,mm/s²). Der Faktor \(\sqrt{\text{rate}/2}\) kommt
        daher, dass die Dichte über die äquivalente Rauschbandbreite integriert wird — bei
        Abtastrate \(f_s\) ist das \(f_s/2\). Fester Bias \(b=0{,}05\,\mathrm{m/s^2}\):
        nach \(T=10\)\,s ist \(\Delta v=bT=0{,}5\)\,m/s und der Wegfehler
        \(\tfrac{1}{2}bT^2=2{,}5\)\,m. Das weiße Rauschen trägt nach derselben Zeit nur
        \(\sigma_a T^2 N^{3/2}/\sqrt{3}\approx 0{,}03\)\,m bei (\(N=1000\) Stiche) — der
        Bias dominiert um zwei Größenordnungen. \emph{Folge:} Die IMU taugt als
        \emph{Bewegungsmodell} mit kurzem Gedächtnis, nicht als Positionsmesser; der Bias
        muss geschätzt (oder zumindest durch Mittelung im Stand klein gemacht) werden.
  \item \textbf{Rekursion.} Mit \(dt=0{,}5\) und \(q=0{,}6\):
        \[
          FP F^{\mathsf T}=\begin{pmatrix}0{,}04+dt^2\cdot1{,}0 & dt\cdot1{,}0\\ dt\cdot1{,}0 & 1{,}0\end{pmatrix}
          =\begin{pmatrix}0{,}29&0{,}50\\0{,}50&1{,}00\end{pmatrix},\quad
          Q=0{,}6\begin{pmatrix}0{,}0417&0{,}125\\0{,}125&0{,}5\end{pmatrix},
        \]
        also \(P^{-}=\begin{pmatrix}0{,}315&0{,}575\\0{,}575&1{,}30\end{pmatrix}\) und
        \(S=0{,}315+0{,}25=0{,}565\). Verstärkung
        \(K=P^{-}H^{\mathsf T}/S=(0{,}315/0{,}565,\;0{,}575/0{,}565)=(0{,}56,\;1{,}02)\).
        Innovation \(y=2{,}40-2{,}00=0{,}40\), damit
        \(\hat x^{+}=2{,}00+0{,}56\cdot0{,}40=2{,}22\)\,m und
        \(\hat v^{+}=0+1{,}02\cdot0{,}40=0{,}41\)\,m/s — die Verstärkung im zweiten Eintrag
        darf größer als 1 sein, sie wirkt auf eine Geschwindigkeit. Position:
        \(\sigma^{-}=\sqrt{0{,}315}=0{,}56\,\mathrm{m}\) vor,
        \(\sigma^{+}=\sqrt{0{,}315-0{,}56^2\cdot0{,}565}=0{,}37\,\mathrm{m}\) nach dem Update.
  \item \textbf{Odometrie-Drift.} Der Posenfehler der Odometrie ist die Summe aller bisher
        begangenen Fehler: er ist zeitlich korreliert, wächst mit der zurückgelegten Strecke
        und hat einen Vorzeichen-zusammenhang (ein zu großer Radradius bleibt zu groß). Weißes
        Rauschen mit festem \(R\) beschreibt aber einen Fehler, der mit jeder Messung
        vergessen wird. Folge: Der Filter hält jede neue Odometrie-Pose für eine unabhängige
        Messung, \(P\) schrumpft gegen null, die Kovarianz ist zu klein — das NEES steigt
        weit über 1, und im GPS-Funkloch driftet die Schätzung unhinterfragt mit. Richtig:
        Odometrie als \emph{Eingang} (Geschwindigkeit) in die Prädiktion benutzen und den
        Rest als Prozessrausch \(Q\) (wachsend mit \(dt\)) oder als zusätzlichen Bias-Zustand
        modellieren.
  \item \textbf{NEES.} \(0{,}1\): Die angegebene Standardabweichung ist im Mittel
        \(\sqrt{1/0{,}1}\approx3{,}2\)-mal zu groß — der Filter ist unnötig vorsichtig und
        traut dem Messwert zu wenig. \(1\): konsistent, die Streuung der Fehler passt zur
        Angabe. \(8\): zu klein um \(\sqrt{8}\approx2{,}8\)-mal — übermütig, und damit
        gefährlich, weil die Anwendung (Regelung, Sicherheit) sich auf die Angabe verlässt.
        Durch das Band \([0{,}15;\,3{,}5]\) für K3 fallen \(0{,}1\) und \(8\), nur \(1\)
        besteht. (Der untere Rand ist absichtlich locker — \(0{,}15\) entsprechen einer
        \(2{,}6\)-mal zu großen Angabe: lässig, aber nicht gefährlich. Über \(3{,}5\) wird es
        eng, weil selbstsichere Falschangaben im Labor niemandem helfen.)
\end{enumerate}

\section{Themen- und API-Übersicht}\label{sec:api}
Der Zugang zu Simulation und Themen läuft über \ordner{mecanum\_lab/robot\_io.py}; unter
ROS 2 wie im Stub-Lauf gleich. Für Versuch 2 sind diese Methoden zuständig:

\begin{table}[ht]\centering\small
\caption{Studentische API in Versuch 2 (Auszug aus \ordner{robot\_io.py}).}
\label{tab:api}
\begin{tabular}{@{}lp{6.0cm}p{3.4cm}@{}}
\toprule
Aufruf & Inhalt & Brauchen Sie für \\
\midrule
\kommando{rob.gps()} & \thema{Gps(t, x, y, theta)} — langsam, verstreut, im Funkloch \kommando{None} & K1--K4 \\
\kommando{rob.odom()} & \thema{Odom(t, x, y, theta, vx, vy, omega)} — 50 Hz, driftend & K2, K4 \\
\kommando{rob.imu()} & \thema{Imu(t, ax, ay, az, gx, gy, gz, roll, pitch)} im Körperframe & K2, K4 \\
\kommando{rob.scan()} & Laser-Entfernungen (Versuch 1, hier nicht bewertet) & — \\
\kommando{rob.age("gps")} & Alter der letzten Messung in s — \kommando{1e9}: nie eine gehabt & überall \\
\kommando{rob.sensor\_profil()} & aktives Prüfprofil als dict — hier steht \(\sigma\) für \(R\) & überall \\
\kommando{rob.config("gps.sigma\_xy", 0.5)} & einzelner Konfigurationswert, mit Standardwert & \(R\), \(Q\) \\
\kommando{rob.task()} & Name des laufenden Auftrags — ändert sich je Teilaufgabe & Loop-Ende \\
\kommando{rob.running()} & läuft die Simulation noch? & Loop-Bedingung \\
\kommando{rob.spin(dt)} & Messungen verarbeiten, ohne die Wanduhr zu befragen & Loop-Körper \\
\kommando{rob.send\_kf(x, y, theta, sx, sy, sth)} & Ihre Schätzung samt \(1\sigma\) — die Bewertungsgröße & K1--K4 \\
\kommando{rob.kf()} & was zuletzt auf \thema{kf/pose}stand (Kontrolle des eigenen Knotens) & Debug \\
\kommando{rob.truth()} & exakte Pose — \emph{im Filterpfad tabu}, für Tests erlaubt & Auswertung \\
\bottomrule
\end{tabular}
\end{table}

Dazu die Themen aus \cref{tab:themen} und, für Neugierige, der Blick von außen:
\kommando{ros2 topic echo /alice/imu --once}, \kommando{ros2 topic hz /alice/kf/pose},
\kommando{ros2 topic echo /sim/config --once}. Wer ohne ROS arbeitet, nutzt dieselben
Themen über den In-Prozess-Bus — nur \kommando{ros2} gibt es dort nicht.

\end{document}
